混入率が低い条件ではStorey型の閾値が$\tau$に達せず，上限付きと上限なしの値は一致する．混入率30\%以上の\NCaprecNHi{}条件（SSB-hard 980クラス$\times$20，50枚）では上限付きのStorey型は$\pP$だけより\NCaprecStoreycHiMaxMinusLpAbs{}--\NCaprecStoreycHiMinMinusLpAbs{}ポイント良く，980クラス$\times$50枚では種に混じるIDの割合が\NCaprecSsbHardCNineEightZeroRFiveZeroStoreyIdfrac{}から\NCaprecSsbHardCNineEightZeroRFiveZeroStoreycIdfrac{}に下がる．全\NCaprecNCells{}条件で上限付きのStorey型が$\pP$だけより1ポイント以上悪いのは\NCaprecStoreycAllNWorseLpOne{}条件（NINCO $c{=}10,20$の$r{=}10$；上限なしも同じ値）で，そこでは種が数十枚しか選ばれず，その\NCaprecNincoCOneZeroROneZeroStoreycIdfracPct\%がIDである（未知がごく少なく$\hat\pi_0$が小さく出る条件；単純BH法は種を選ばない）．単純BH法は上限の有無に依らず全条件で$\pP$だけより最大\NCaprecBhcAllMaxMinusLp{}ポイントしか悪くならない．登録時の仮説H3（付録\ref{sec:app_hyp}：混入率30\%以上で$\pP$だけより1ポイント以上悪くならず，$r\ge10$・混入率30\%未満で上限なしの利得を1ポイント以内に保つ）は，後者の\NCaprecLoROneZeroN{}条件での差が最大\NCaprecLoROneZeroMaxStoreycMinusStorey{}ポイントであり，成り立つ．
